\section{Introduction}\label{sec:introduction}

Canard dynamics turn a small parameter displacement into a large change of
orbit geometry.  In ordinary slow--fast systems this mechanism is organized
locally near a loss of normal hyperbolicity
\citep{krupa2001extending,krupa2001relaxation,wechselberger2012canards}.
Delay equations add a more basic ambiguity: the state is a history, and a
current-state projection can discard precisely the transverse delay response
that later returns to the critical direction.  Delayed canards and delayed
FitzHugh--Nagumo oscillations are by now well established
\citep{campbell2009delay,krupa2016complex,krupa2016explosion}, but a local
canard calculation does not automatically survive embedding into a large
retarded network.

A second threshold problem appears when the system is actuated.  A finite
physical pulse generates a curve in the RFDE history space.  A selected
pulse/stable-sheet threshold is a value at which that curve meets a declared
codimension-one stable sheet.  It becomes a biological onset only after the
two sides are proved to route to different attracting objects.  Neither
notion is the same as a Lin-gap canard root, and a projected observable event
is not by itself any of the three.  Our organizing principle is therefore to
keep these objects separate while exposing the transverse mechanism common
to them.

\subsection{Three levels of generality}

The paper distinguishes three statements.

First, canonical synchrony gives an exact upper-triangular block
decomposition, and the transverse complete inverse gives a bounded codomain
row reduction to the scalar/transverse direct sum.  When all endpoint, phase,
parameter, and gap data are collective, the corrected full cokernel
transfers the scalar Fredholm data and gap without changing the root.

Second, a one-gap structural theorem allows small network residuals that need
not preserve exact synchrony.  Its constants are independent of the number
of nodes, provided the critical projections, transverse estimates, history
traces, model-fitting map, and simple-root bounds are uniform.  This is a
local selected-root theorem, not a claim about arbitrary graph sequences.

Third, a shared-resource heterogeneous-curvature class realizes the
mechanism on finite directed Markov networks with a common Dobrushin gap and
without a nontrivial synchrony quotient.  Fixed delay support and a declared
row-neutral perturbation are essential.  Sparse or closing-gap families and
arbitrary changes of the base topology are outside the proved scope.

\subsection{Contributions and claim boundary}

The theorem-level contributions currently assembled in this draft are:

\begin{enumerate}[label=(\roman*)]
\item a dimension- and topology-uniform transverse decay estimate and causal
  complete-line Green bound for every finite leaky Dobrushin network with a
  nonnegative row-stochastic current layer and fixed nonnegative half-row-mass
  delay layers in the declared voltage strip, without common delay-layer
  left balance;
\item an exact upper-triangular canonical Lin factorization, bounded
  invertible codomain row reduction to the scalar/transverse direct sum, and
  corrected full cokernel, together with a general normalized one-gap
  perturbation theorem having an explicit root radius and quadratic remainder;
\item a synchrony-quotient-free shared-resource network whose selected
  complete-history root has a nonzero, dimension-uniform topology-resolvent
  coefficient;
\item validated scalar leaky-recovery invariant objects: a stable quiet
  equilibrium, two periodic RFDE branches, exactly one nontranslation
  unstable multiplier for the center inner orbit, and a directed nonsingular
  frequency--amplitude response, together with a continuous RFDE adjoint and
  a correlated physical-time base-orbit \(uu\) stable-output bound, a direct
  common-row terminal certificate for the selected phase-fixed inner linear
  map with \(\|AP_s\|_{\Sigma_0\to Y}\le0.009896427481610001<0.1\), yielding
  \(K_s=1\) at the registered rate \(\rho_s=0.1\), and an event-aware
  signed-density proof that the phase-fixed outer linear return contracts
  arbitrary reduced \(C^0\) section histories with norm at most \(0.550516\),
  followed by a source-bound nonlinear local return and attracting tube with
  section radius \(r_6=10^{-335}\), full complete-history tube radius
  \(R_6=10^{-3}\), and return constant \(\Lambda_6\le0.561839\);
\item a general eventual-smoothing theorem for selected complete-history
  RFDE events, together with a model-specific common near-two-period event
  tube on the explicit arbitrary-\(C^0\)-history scale
  \(\lambda_0=1.2\times10^{-8}\), obtained from an orbit-aware weighted
  energy estimate, the direct derivative identity
  \(DQ_Y(Y_*)=A^2\), transferred two-step hyperbolic rates, and a finite-
  sampling obstruction showing why nodal refinement alone cannot certify
  the required Banach-space Hessian;
\item a wide-parameter physical-pulse theorem with a unique transverse
  Route--C event, a fourth-order event-time graph, and a continuous
  complete-history tube on the whole parameter interval, together with a
  direct event-aligned first-variation theorem retaining the moving-time
  contribution and a fixed-center, physical-right-gauge-corrected normalized
  action that has one strict sign throughout the pulse interval, together
  with source-bound endpoint functional signs, a full-interval
  stable-coordinate ball, and a strengthened conditional stable-gap slope
  gate in the same section, history norm, and Grushin normalization;
\item an event-aligned stable-sheet theorem and a block-triangular
  frequency--amplitude--threshold-offset inverse, conditional on explicit
  numerical graph and crossing inequalities; its interpretation as a safety
  coordinate additionally requires the stated routing inequalities.
\end{enumerate}

Item (vii) is not promoted to a model theorem in the present research draft.
The discrete stable-power input required by a future inner graph transform is
now proved for the selected near-one-period phase-fixed linear map, with
\(K_s=1\) and \(\rho_s=0.1\).  The selected two-period map is now defined and
\(C^2\) on a microscopic nonzero scaled ball, and its center derivative,
fixed splitting, and two-step rates are proved.  The remaining inner
certificates are its enlargement to the preferred graph domain, return-domain
containment there, all six uniform split-ball Hessian blocks, and a
quantitative stable graph whose domain contains the certified
radius-\(47/5000\) stable-coordinate ball,
together with its derivative and endpoint-height bounds.  The stable-coordinate
ball itself and the conditional strict-slope implication are proved;
graph-adjusted endpoint signs and the selected crossing are not.  Entry of the
\(J=0.32\) history into the microscopic outer return domain, outer-side
attachment, and two-sided basin routing also remain open.  The proved Route--C
event is a common history-space section, not an ordinal claim about the third
post-release crossing.  A global outer Floquet count is useful but is not
logically required if local attraction is proved directly.  Conversely, none
of the network root results identifies the selected canard root with the
physical pulse onset.

Frequency--amplitude reduction and control of biological oscillators motivate
the proposed final coordinate system
\citep{wilson2016isostable,kotani2020nonlinear,qin2021coordinator,
zhong2023modulating}.  Its third target coordinate is the signed offset from
a selected history-space stable-sheet threshold, not an independently
imposed detector variable.  Only the still-open uniform routing theorem would
promote that offset to a biological safety margin.

\subsection{Organization}

\Cref{sec:model-results} introduces the scalar RFDE, finite network class,
and theorem package.  Subsequent sections will supply the complete-line
inverse, structural root transfer, validated scalar objects, physical-pulse
separator, and network robustness.  Directed arithmetic, source manifests,
and hostile mutation tests belong to the reproducibility appendices rather
than to the mechanism-level argument.
